% Generated draft: requires mathematical and source review.
\documentclass[11pt]{article}
\usepackage[margin=25mm]{geometry}
\usepackage{fontspec}
\setmainfont{Latin Modern Roman}
\usepackage{amsmath,amssymb,mathtools,unicode-math}
\setmathfont{Latin Modern Math}
\usepackage{enumitem,needspace,xurl,hyperref}
\setlength{\parindent}{0pt}
\setlength{\parskip}{5pt}
\setlength{\emergencystretch}{3em}
\newcommand{\N}{\mathbb{N}}
\newcommand{\Z}{\mathbb{Z}}
\newcommand{\Q}{\mathbb{Q}}
\newcommand{\R}{\mathbb{R}}
\newcommand{\C}{\mathbb{C}}
\newcommand{\F}{\mathbb{F}}
\newcommand{\A}{\mathbb{A}}
\newcommand{\PP}{\mathbb{P}}
\newcommand{\E}{\mathbb{E}}
\newcommand{\eps}{\varepsilon}
\newcommand{\dd}{\,\mathrm{d}}
\begin{document}
\section*{UnsolvedMath ID: 30005013}
\textbf{Source problem code:} catalogue entry
\subsection{Definitions and mathematical statement}
For an integer \(k\geq 0\), let
\[
\Pi_k^3=\operatorname{span}\{x_1^{a_1}x_2^{a_2}x_3^{a_3}:a_i\in\Z_{\geq0},\ a_1+a_2+a_3\leq k\},
\qquad \dim\Pi_k^3=\binom{k+3}{3}.
\]
For a domain \(\Omega\subseteq\R^3\) and an integrable nonnegative weight \(w\) with the required finite moments, a cubature formula exact through total degree \(m\) is an identity
\[
\int_\Omega p(x)w(x)\,\dd x
=\sum_{j=1}^{N}\lambda_jp(\xi_j)
\qquad\text{for every }p\in\Pi_m^3,
\]
where the nodes \(\xi_j\in\Omega\) are distinct and the coefficients \(\lambda_j\) are real. Positivity of these coefficients is an additional condition whose intended use must be confirmed from the source.

The supplied question asks whether a Gaussian cubature formula exists for a three-variable product Chebyshev weight. Its prescribed degree and precise weight are missing from the clean statement. The excerpt discusses both \(m=2n-1\) and \(m=2n-2\). A candidate formulation, conditional on the intended convention, asks for the identity above with
\[
N=\dim\Pi_{n-1}^3=\binom{n+2}{3}
\]
for a specified integer \(n\geq1\), with \(m\) chosen from those two degrees. The record does not specify which degree, which values of \(n\), or whether existence is requested for one value, every value, or an infinite family of values.

If the intended domain is the three-dimensional analogue of the displayed square, the two conventional product Chebyshev candidates on \([-1,1]^3\) are
\[
w_{\mathrm{I}}(x)=\prod_{i=1}^{3}(1-x_i^2)^{-1/2},
\qquad
w_{\mathrm{II}}(x)=\prod_{i=1}^{3}(1-x_i^2)^{1/2}.
\]
These formulas specify densities on the interior; integrals over the closed cube are interpreted using these integrable densities, without assigning finite endpoint values to \(w_{\mathrm{I}}\). They are alternatives for clarification, not a selection of the intended weight. The excerpt calls its two-variable weight ``of the second kind,'' but its damaged displayed formula appears to contain reciprocal square roots. It does not explicitly supply a three-variable density, normalization, or boundary-node convention.

Consequently, the supplied evidence does not determine a single precise existence proposition. The intended degree, parameter quantifiers, weight, domain, node restrictions, coefficient restrictions, and meaning of ``Gaussian'' must be resolved before the conditional formulation can be accepted as the mathematical statement.
\subsection{Short English statement}
The question concerns an exact integration formula for polynomials in three variables using a Gaussian number of nodes and a product Chebyshev weight. The supplied record leaves the degree and parameter quantifiers unspecified, and its weight name conflicts with its damaged formula. The displayed formulation therefore records conditional interpretations requiring human mathematical review.
\subsection{Sources}
\begin{itemize}
\item Local registry: supplied UnsolvedMath source record, permanent ID \(30005013\), source problem code \texttt{catalogue entry}, Oberwolfach Reports Open Problems collection. No proposer is identified; the supplied proposed year is 2022.
\item Supplied report citation: \emph{Mini-Workshop: Interpolation, Approximation, and Algebra} (2023), Oberwolfach Reports, DOI 10.4171/owr/2022/7, pp. 383--428. \url{https://doi.org/10.4171/owr/2022/7}. The supplied background describes these reports as records of mathematical workshops and open-problem discussions. This source has not been independently consulted for this draft.
\item Supplied literature source: De Marchi--Vianello--Xu, cube hyperinterpolation. \url{https://www.sciencedirect.com/science/article/pii/S0898122107007493}.
\item Supplied literature source: Li--Sun--Xu, cubic-domain cubature. \url{https://www.global-sci.com/nmtma/article/view/13901}.
\item Literature assessment dated exactly 2026-08-22: the supplied curation labels the problem open and states that three-dimensional product-Chebyshev formulas with the required exactness and near-minimal node counts are known, while a minimal or Gaussian rule of the requested type had not been established. This is supplied dated curation derived from the OWR material, not an independently verified literature assessment; its status and cited claims require independent verification.
\item The registry labels the statement \texttt{corrected\_verified} and reports that it was checked against the official Oberwolfach report. This is a supplied verification claim, not verification performed for this draft. The excerpt also refers to ``[1, Chapter 3],'' but the corresponding bibliographic entry is absent.
\end{itemize}
\end{document}
